% Cover letter for the Fluid Phase Equilibria submission; built to builds/cover_letter.pdf by scripts/build_manuscript.sh.
\documentclass[11pt]{article}
\usepackage[letterpaper,hmargin=0.8in,vmargin=0.65in]{geometry}
\usepackage[T1]{fontenc}
\usepackage{newtxtext,newtxmath}
\usepackage[version=4]{mhchem}
\usepackage{xspace}
\usepackage{siunitx}
\usepackage{enumitem}
\usepackage{xurl}
\usepackage[hidelinks]{hyperref}
\urlstyle{same}
\input{chemistry_macros}
\setlength{\parindent}{0pt}
\setlength{\parskip}{0.4\baselineskip}
\pagestyle{empty}
\setlist[itemize]{leftmargin=1.2em,itemsep=0.2\baselineskip,topsep=0.1\baselineskip}

\begin{document}

\begin{center}
{\large\textbf{Tanner W. Polley}}\\[2pt]
Department of Chemical Engineering, Brigham Young University, Provo, UT 84602, United States\\
\href{mailto:tpolley3@byu.edu}{tpolley3@byu.edu}
\end{center}
\vspace{-0.6\baselineskip}
\rule{\linewidth}{0.4pt}

\today

The Editors, \textit{Fluid Phase Equilibria}

Dear Editors of \textit{Fluid Phase Equilibria},

Please consider ``Modeling \COtwo solubility in aqueous alkanolamine solutions with ePC-SAFT'' as a research article.

The paper describes \COtwo solubility in four aqueous amines with one ePC-SAFT model with explicit ions.
The four calibrations share one base set that no amine data determine: water, \COtwo, the inorganic ions and reactions, and water--anion interactions from osmotic coefficients of the sodium salts.
Monoethanolamine (MEA) is fitted jointly to pressures, species fractions, absorption heats, heat capacities and densities.
Diethanolamine (DEA), methyldiethanolamine (MDEA) and 2-amino-2-methyl-1-propanol (AMP) are then fitted one by one, each through its own species, and every fitted model is searched for a second liquid.

What is new:
\begin{itemize}
\item \textbf{Four amines on one base without amine data.}
The fitted \COtwo partial pressures of MEA at 30\,wt\% and of DEA deviate by \res{mea-fitted-pressure-aard-abstract} and \res{dea-fitted-pressure-aard-abstract}\,\%, and the fitted total pressures of MDEA and AMP by \res{mdea-fitted-pressure-aard-abstract} and \res{amp-fitted-pressure-aard-abstract}\,\%.
DEA and AMP each hold two parameters at their MEA values, so the base and these parameters give the pressures of a primary, a secondary, a tertiary and a hindered amine from one Helmholtz energy.
\item \textbf{Water--anion interactions from salt data.}
The water--bicarbonate and water--carbonate interactions are fitted to osmotic coefficients of the sodium salts and adopted because data without an amine determine them.
The supplement compares the deviations with those of an interaction fitted to the MEA data.
\item \textbf{Phase stability of the fitted models.}
Reactive, electroneutral tangent-plane searches at measured and gridded states find no confirmed second liquid in the four models, and the fitted data decide that outcome.
A DEA fit without the Lawson--Garst pressures predicts a second liquid.
Weighting the \COtwo-free MDEA pressures like the loaded ones removes a calculated liquid--liquid gap from that solution.
An AMP fit without the loaded Tong pressures and the Hartono pressures above \SI{353}{\kelvin} predicts a second liquid, and the AMP model, which fits them, predicts none.
A phase-stability search therefore tests a fitted electrolyte model in a way that its deviations from the fitted data do not.
\end{itemize}

The work fits the journal's scope in equations of state for electrolyte and reacting systems, chemical and phase equilibrium, and phase-stability analysis.
It builds on ePC-SAFT studies of aqueous amines and electrolytes published in \textit{Fluid Phase Equilibria}.
The MEA model is compared with the author's earlier reactive ePC-SAFT model of MEA (Polley and Hedengren, \textit{Next Chemical Engineering} 3 (2026) 100132), whose parameters were estimated in separate steps.

The four models supply thermodynamic properties for absorber design within the ranges and deviations that the limitations section states property by property.
That section also gives the conditions where the deviations are larger, among them MEA pressures at 15 and 45\,wt\% and absorption heats at \SI{393}{\kelvin}.

A supplement gives the model contributions, parameters, data conventions, deviations by data set and phase-stability details.
The electronic data comprise processed observations, parameter values, numerical results and the analysis and figure code; the deposit also contains the epcsaft build used for the reported values, whose source is public.
The deposit is archived on Zenodo at \url{https://doi.org/10.5281/zenodo.24623748}.
The manuscript includes funding, competing-interest, CRediT and generative-AI declarations.

Sincerely,\\[1.4\baselineskip]
Tanner W. Polley

\end{document}
